% Теоретические аспекты цикла
\section{Theoretical aspects of the cycle}

\begin{frame}
  \frametitle{Eclectic cycle definition}
  The multitude of competing theoretical and methodological approaches to analyzing the cycle in the economy necessitates the formulation of an eclectic definition as the initial abstraction for further theoretical specification of the concept:

  \begin{block}{Definition}
    \textbf{Economic cycle} – оobjectified wave–like movement of aggregate economic activity within the framework of patterns of changes in aggregate output and other related indicators of the economic environment (production, employment, prices, profits, investments etc.), described as a sequential change in the phases of its qualitative states and externally observed in the form of stable deviations from a certain \textquotedbl{}normal\textquotedbl{} state or a long-term development trend.
  \end{block}

  \begin{alertblock}{Important}
    This definition is extremely abstract and mainly descriptive, therefore is not fully reliable and requires further specification.
  \end{alertblock}
\end{frame}

\begin{frame}
  \frametitle{Methodology of cycle theories}
  All existing cycle theories can be classified into the following categories based on their methodology:
  \begin{itemize}
    \item \textit{marxist theories:} consider the industrial cycle as a specific form of dialectical development of an internally contradictory capitalist economy within a certain historical stage;\footnote{Marx, Engels, Maksakovskiy, Mendelson et al.}
    \item \textit{bourgeois theories:} consider the phenomenon of the cycle through the lens of the absolutization of the role of individual factors within an internally consistent economic system;\footnote{Kitchin, Kuznets, Juglar, Kydland \& Prescott, Schumpeter et al.}
    \item \textit{eclectic theories:} combine incompatible elements of the two previous categories in various proportions or externally reproduce the conceptual framework of one of them while actually borrowing the substantive provisions of the other.\footnote{Kondratiev, Luxembourg, Bukharin et al.}
  \end{itemize}
\end{frame}

\begin{frame}
  \frametitle{Criteria for analysis and systematization}
  In order to further systematize scientific works and specify their theoretical provisions, a number of key criteria have been identified:

  \begin{columns}[T,onlytextwidth]
  \begin{column}{0.48\textwidth}
    \begin{enumerate}
      \item Definition
      \item Terms
      \item Methodology of research
      \item Scientific approaches
      \item Tools of research
      \item Methods of research
    \end{enumerate}
  \end{column}

  \begin{column}{0.48\textwidth}
    \begin{enumerate}
      \setcounter{enumi}{6}
      \item Indicators
      \item Economy
      \item Initial sphere
      \item Fundamental causes
      \item Factors
      \item Period
      \item Multicyclicity
      \item Phases
      \item Regulation
    \end{enumerate}
  \end{column}
\end{columns}
\end{frame}

\begin{frame}
  \frametitle{Industrial cycle (1)}
  \begin{alertblock}{Finding}
    A critical analysis of existing theories reveals that the Marxist theory of the industrial cycle provides the most comprehensive and accurate description of cyclical phenomena within a capitalist economy.
  \end{alertblock}

  \begin{block}{Definition}
    \textbf{Industrial cycle} – objectively determined, historically specific and relatively periodic form of the movement of expanded reproduction of total social capital, expressed in the regular and sequential alternation of phases of crisis, depression, recovery, and growth, arising from the accumulation, exacerbation, temporary violent resolution, and subsequent reproduction of the internal contradictions of the capitalist mode of production.
  \end{block}
\end{frame}

\begin{frame}
  \frametitle{Industrial cycle (2)}
  \begin{tabularx}{\textwidth}{l>{\raggedright\arraybackslash\hyphenpenalty=10000\exhyphenpenalty=10000}X}
    \textbf{Indicators} & \textit{Industrial production volume} – general dynamics and phases; \textit{GDP} – initial point of general overproduction crisis \\
    \textbf{Economy} & Developed capitalist economy with a predominant share of large-scale machine production; national economy; 1825 – the first crisis of general overproduction in Great Britain \\
    \textbf{Initial sphere} & Industrial production \\
    \textbf{Fundamental causes} & Immanent contradictions of the capitalist system \\
    \textbf{Factors} & Massive periodic renewal of fixed capital
  \end{tabularx}
\end{frame}

\begin{frame}
  \frametitle{Industrial cycle (3)}
  \begin{tabularx}{\textwidth}{l>{\raggedright\arraybackslash\hyphenpenalty=10000\exhyphenpenalty=10000}X}
    \textbf{Period} & $\approx$ 10 years, due to the period of physical and moral depreciation of fixed capital \\
    \textbf{Multicyclicity} & The dominant role of the industrial cycle, partial impact on the dynamics of other cycles\tablefootnotetext{In particular, the agricultural cycle} \\
    \textbf{Phases} & Four-phase model: \newline crisis $\to$ depression $\to$ recovery $\to$ growth \\
    \textbf{Regulation} & Government regulation can accelerate or delay the onset of a crisis, change its nature or duration, but cannot eliminate it
  \end{tabularx}
\end{frame}